\chapter{Results and Analysis} \label{Chap4}\label{chap:results}

\section{Evaluation scope}

Results are organised around the evaluation questions: data integrity,
large-image density and colour inference, independent-mesh behaviour,
coordinate control, static and on-demand integration, matched representation,
and component verification.  Quantitative findings are tied to the data population or
processing stage from which they were measured.

\section{Prototype feasibility and empirical focus}

A retained London GLB and corresponding server record establish compatibility
with the released Sat3DGen inference interface.  The Unity/Ubiq client also
compiled with the intended acquisition, request, retrieval, and scene-loading
stages, informing the design of the downstream interface.  The principal
experiments comprise a fresh contiguous London large-image run, a matched
overlap-fusion comparison, global vertex-colour recovery, an independent-OBJ
mesh-processing study, and ChordAtlas publication.  Retained ChordAtlas
appearance exports provide representation-level evidence, while the earlier
format-conversion prototypes contribute architectural context.

\section{Data construction and integrity}\label{sec:data-results}

The London builder is assessed for dataset readiness and downstream geospatial
reuse.  The results address coverage, independence, label geometry, and
auxiliary-channel integrity rather than model-training performance.

\subsection{Inventory and physical extent}

Panorama acquisition covered the majority of target locations, while the satellite, DSM, OSM, and grid channels cover only a subset; sampling density and overlap also differ between branches.

Projected tile unions expose the difference between file count and physical
coverage.  The 127-tile set covers approximately
\SI{0.45}{\kilo\metre\squared}; the repaired grid covers
\SI{0.96}{\kilo\metre\squared}.  The 50 and 10 per cent overlap layouts cover
approximately \SI{0.93}{\kilo\metre\squared} and
\SI{0.89}{\kilo\metre\squared}, respectively.  Despite containing 306 tiles,
the latter supports 2,331/2,333 queries because its stride and boundary handling
omit two locations.  Tile count is therefore an insufficient coverage measure.

\begin{figure}[H]
\centering
\includegraphics[width=\textwidth]{data_builder_coverage_channels.pdf}
\caption{Estimated grid-footprint union, matched-query coverage, and stored
channel availability.  Areas use a local-scale Web-Mercator estimate.}
\label{fig:data-builder-coverage-channels}
\end{figure}

\subsection{DSM channel}

The four Environment Agency inputs are adjacent 5,000-by-5,000 float rasters at
\SI{1}{\metre} resolution in EPSG:27700.  The product is the Composite 2022 DSM
\cite{environmentAgencyDSM2022}; it measures visible surfaces, not bare earth.
The source rasters cover all 2,343 query points and all 272 repaired
grid centres.

Figure~\ref{fig:data-builder-dsm-views} visualises the metric source mosaic
before per-patch percentile normalisation.  Mosaicking the four rasters across
the approximately 1.00-by-0.94-kilometre footprint provides complete coverage
at the displayed sampling grid.

\begin{figure}[H]
\centering
\includegraphics[width=\textwidth]{data_builder_dsm_views.pdf}
\caption{North-up and oblique views of the Environment Agency 2022
\SI{1}{\metre} first-return DSM mosaic over the London study area.  The plan
view preserves equal east--north scales and NoData; the oblique view uses
threefold vertical exaggeration.  Colour shows elevation relative to Ordnance
Datum Newlyn.  Source: Environment Agency, LIDAR Composite DSM 2022; Open
Government Licence v3.0
\cite{environmentAgencyDSM2022,nationalArchivesOGL3}.}
\label{fig:data-builder-dsm-views}
\end{figure}

Independent 5th--95th-percentile scaling makes the patches suitable as relative
surface-height channels but removes cross-patch comparability in absolute
elevation.  A legacy DSM derivative without traceable production metadata was
therefore excluded from the evaluated channel.

\subsection{OSM semantic channel}

The Overpass query retrieves support nodes recursively, and the extractor emits
all returned nodes as point features even when they have no category tags.  Of
67,281 category records, 57,623 are such support-node records (85.65 per cent).
Multipolygon fragments are also treated as
independent outer rings, and inner-ring assignment can fill intended holes.
Semantic validity therefore requires geometry- and tag-aware validation rather
than file-count checks.
The active mesh loader adds 19,682 ordinary-building and 9,237 height-building
records to the same index
\cite{openstreetmapData}.

\section{Large-image density and colour inference}
\label{sec:large-image-results}

\subsection{Contiguous London configuration}

The contiguous run covered approximately \SI{285.5}{\metre} by
\SI{475.8}{\metre} from a 3,072-by-5,120-pixel zoom-20 mosaic.  Each
640-pixel source window was passed through the unchanged 256-pixel Sat3DGen
image transform.  A 160-pixel stride produced a 17-by-29 schedule, or 493
windows, with 75 per cent nominal overlap.  The mesh-resolution setting was
192 voxels; it did not change the network image size.  After fractional
placement, the fused density volume measured $1232\times740\times173$ and had
no zero-weight cells.

One global Marching Cubes extraction produced 3,122,922 vertices and 6,241,862
faces.  Density inference and fusion required 129.0 seconds and reached
\SI{6.62}{\gibi\byte} peak allocated GPU memory.

\subsection{Window-boundary comparison}

The hard-box and fractional-feather runs used the same source image,
checkpoint, windows, overlap, voxel resolution, and isovalue.  The baseline is
described as release-style because it retains the released wrapper's hard
boundary crop and equal-weight overlap rule, while using the matched project
settings rather than the release defaults.  Figure~\ref{fig:large-image-feather}
shows the resulting top surfaces.

\begin{figure}[H]
\centering
\includegraphics[width=\textwidth]{big_image_feather_comparison.pdf}
\caption[Hard-box and fractional-feather large-image fusion.]{Matched
large-image fusion comparison.  Hard-box averaging produces repeated
window-aligned height steps; fractional placement and raised-cosine feathering
replace the abrupt contributor changes with smooth overlap transitions.  The
false colour and shading encode density-derived top-surface height, not
satellite texture or appearance accuracy.}
\label{fig:large-image-feather}
\end{figure}

At the former hard-crop contribution boundaries, mean absolute density jumps
fell from 0.710 to 0.286 along rows and from 0.963 to 0.491 along columns
(Table~\ref{tab:large-image-seams}).  Mean top-surface jumps fell from 2.596 to
0.659 voxels and from 3.065 to 1.214 voxels, respectively.  The corresponding
reductions are 49--60 per cent for density and 60--75 per cent for the derived
top surface.  Geometry runtime increased by 8.6 per cent, while peak allocated
GPU memory was unchanged at the reported precision.  The comparison therefore
supports a reduction in boundary discontinuity.

\begin{table}[H]
\centering
\caption{Matched hard-box and fractional-feather fusion results.}
\label{tab:large-image-seams}
\small
\begin{tabularx}{\textwidth}{Xrrr}
\toprule
Measure & Hard-box & Fractional feather & Change\\
\midrule
Mean density jump, row boundary & 0.710 & 0.286 & $-59.8\%$\\
Mean density jump, column boundary & 0.963 & 0.491 & $-49.0\%$\\
Mean top-surface jump, row (voxels) & 2.596 & 0.659 & $-74.6\%$\\
Mean top-surface jump, column (voxels) & 3.065 & 1.214 & $-60.4\%$\\
Geometry stage (s) & 118.7 & 129.0 & $+8.6\%$\\
Peak allocated GPU memory (GiB) & 6.62 & 6.62 & unchanged\\
\bottomrule
\end{tabularx}
\end{table}

\subsection{Global vertex-colour recovery}

The second pass accumulated 42.6 million positive window contributions over
the fixed global surface.  Every vertex received support from at least one
window, with no uncovered or all-black vertices.  Reloading the coloured PLY
preserved every vertex position and face index and reproduced the exported RGB
values.  The colour stage required 97.3 seconds, of which 80.9 seconds were
model queries, and reached \SI{1.86}{\gibi\byte} peak allocated GPU memory.
The exported PLY displayed its blended vertex colours in MeshLab.  These checks
establish complete and index-consistent appearance transfer.

\subsection{Current myProject large-image publication}

The large-image path was also exercised through the ChordAtlas bridge on a
1,536-by-1,536-pixel, 49-window selection.  Its retained job records a
340,825-vertex, 678,726-face density surface followed by complete colour
support.  The geometry stage required 17.0 seconds and the colour stage 23.2
seconds.  The final colour-enabled publication reused the validated geometry
stage from the same retained job before running and validating the second pass.

After one local-coordinate conversion, the context crop contained 66,666
vertices and 132,440 faces.  DSM processing covered every retained vertex and
preserved its RGB index while changing vertical position by at most
\SI{30.7}{\metre}.  Footprint extraction published two non-exclusive coarse
building assets.  Their geometric union, represented by the top-level
selection crop, contains 17,571 vertices and 34,383 faces; the two building
files share 266 vertex positions and 452 triangles.  Both OBJ files contain
RGB on every vertex and pass the Java colour parser; per-building MiniMesh
trees also exist.  In ChordAtlas, the RGB OBJ is displayed by a dedicated
vertex-colour renderer, while MiniMesh remains a separate geometric workspace
representation.

\begin{figure}[H]
\centering
\includegraphics[width=\textwidth]{big_image_coloured_publication.pdf}
\caption{Oblique and plan views of the current large-image building
publication after coordinate conversion, DSM processing, and footprint
extraction.  Face colour is the mean of the predicted per-vertex RGB values;
the figure demonstrates colour retention through publication, not comparison
with a reference texture.}
\label{fig:large-image-coloured-publication}
\end{figure}

\section{Mesh-space post-inference experiment}\label{sec:sat-results}

\subsection{Evaluated processing scope}

The experiment processes 15 coordinate-paired OBJ tiles through OSM/DSM
loading, crop and placement, lower-surface removal, stitching, DSM correction,
and building extraction.  The associated cache records describe geometry but
do not bind a model checkpoint and configuration, so the measurements below
characterise the mesh-processing stages.

Qian et al. report FID 19.2 and DSM RMSE \SI{5.20}{\metre} on VIGOR-OOD
\cite{qian2026sat3dgen}.  Those metrics assess the published model protocol;
this experiment measures downstream operations on retained mesh tiles and uses
the external DSM as a processing input.  Unlike the density-domain experiment
in Section~\ref{sec:large-image-results}, it combines independently extracted
meshes and therefore evaluates a complementary assembly strategy.

\subsection{Geographic composition of the mesh batch}

Recursive discovery selected two disconnected clusters: nine eastern tiles in
a 3-by-3 arrangement and six western tiles in a 2-by-3 arrangement.  Their
centres are approximately \SI{666}{\metre} apart, more than eleven nominal tile widths.
Figure~\ref{fig:tile-clusters} plots every centre and nominal footprint.

\begin{figure}[H]
\centering
\includegraphics[width=0.94\textwidth]{tile_clusters.pdf}
\caption{Geographic composition of the 15 OBJ meshes selected by the merge
experiment.  Filename coordinates identify two disconnected sites.}
\label{fig:tile-clusters}
\end{figure}

The merged output consequently combines two sites rather than one contiguous
London scene.  This geographic inconsistency originates in cache selection,
not the upstream Sat3DGen generator \cite{qian2026sat3dgen}.

\subsection{Geometry retained through processing}

Table~\ref{tab:stage-counts} gives the exact geometry counts.  From raw inputs
to the stitched scene, 1,724,759/3,640,214
vertices (47.38 per cent) and 3,184,234/7,281,608 faces (43.73 per cent) were
removed or consolidated.  This large, intended reduction makes the crop,
bottom-removal, and stitching predicates central to output validity.

\begin{table}[H]
\centering
\caption{Geometry counts by post-processing stage.}
\label{tab:stage-counts}
\small
\begin{tabular}{lrrr}
\toprule
Stage & Vertices & Faces & Change from previous stage\\
\midrule
Raw 15 OBJ inputs & 3,640,214 & 7,281,608 & ---\\
Cropped & 2,828,613 & 5,615,039 & $-811{,}601$ V; $-1{,}666{,}569$ F\\
Bottom removed & 2,252,373 & 4,474,289 & $-576{,}240$ V; $-1{,}140{,}750$ F\\
Stitched & 1,915,455 & 4,097,374 & $-336{,}918$ V; $-376{,}915$ F\\
DSM scene & 1,915,455 & 4,097,374 & topology unchanged\\
DSM building extraction & 1,001,507 & 2,116,136 & branch result\\
No-DSM building extraction & 907,370 & 1,911,272 & branch result\\
\bottomrule
\end{tabular}
\end{table}

Figure~\ref{fig:tile09-bottom-removal} shows the saved crop-stage and
bottom-removal outputs for one representative tile.  The horizontal footprint
is unchanged, while the operation removes a near-horizontal generated support
sheet spanning almost the full tile.  Every tile lost exactly 38,416 vertices
and 76,050 faces at this stage, consistent with removal of the same standard
lower grid.

\begin{figure}[H]
\centering
\includegraphics[width=0.98\textwidth]{tile09_bottom_removal.pdf}
\caption[Bottom removal in representative tile 09.]{Representative tile~09 at
two consecutive saved stages.  Panel (a) is the cropped, placed, and
OSM-pre-aligned mesh before bottom removal, with the removed lower sheet
highlighted in orange; panel (b) is the retained mesh immediately afterwards.
Both use the same underside camera and metric frame, and stored vertex colour
is omitted to expose the geometric change.}
\label{fig:tile09-bottom-removal}
\end{figure}

The controlled bottom-removal fixture also showed that the active predicate
removes a triangle when only one vertex is close to the lowest surface,
although the documented rule requires all three vertices to satisfy the
threshold.  In the stitching fixture, two triangles separated by
\SI{10}{\metre} in Y collapsed from six vertices to three because matching used
X/Z coincidence without the configured vertical-separation test.  In the
north-western two-by-three subset used below, the six pre-stitch tiles contain
819,191 vertices and 1,626,029 faces; the saved post-stitch subset contains
696,199 vertices and 1,485,998 faces with the same horizontal extent.  The
sub-metre consolidation is quantified here because it is not distinguishable
in a scene-scale silhouette.

\subsection{DSM displacement is scene-wide}

DSM correction changed all 1,915,455 vertices in the saved scene and every
change was positive.  The median displacement is \SI{36.08}{\metre}, the mean
is \SI{42.24}{\metre}, and values range from \SI{15.81}{\metre} to
\SI{62.63}{\metre}.  The scene's vertical range changes from
\SIrange{-26.77}{18.76}{\metre} to \SIrange{1.86}{67.94}{\metre}
(Figure~\ref{fig:dsm-displacement}).

\begin{figure}[H]
\centering
\includegraphics[width=\textwidth]{dsm_displacement.pdf}
\caption{Per-vertex vertical displacement introduced by DSM correction.  The
distribution characterises implementation behaviour rather than surveyed
height error.}
\label{fig:dsm-displacement}
\end{figure}


\subsection{North-western stitched scene and representative topology}

Figure~\ref{fig:dsm-scene-nw-2x3} isolates the north-western two-by-three
cluster in the stitched scene before and after the saved DSM stage.  The two
OBJ files preserve vertex order, horizontal coordinates, RGB attributes, and
face indices within this subset, so the stage changes only Y.  The subset
contains 696,199 vertices and 1,485,998 faces.  All paired vertices move
upward; median displacement is \SI{38.58}{\metre}, with a 5th--95th percentile
range of \SIrange{32.96}{60.49}{\metre}.

\begin{figure}[H]
\centering
\includegraphics[width=0.98\textwidth]{dsm_scene_nw_2x3_comparison.pdf}
\caption[DSM-stage vertical transformation in the north-western two-by-three
cluster.]{DSM-stage transformation in the north-western two-by-three stitched
scene cluster.  Panel (a) is the post-stitch output before DSM processing;
panel (b) is the same mesh after DSM height correction.  Both use uniform
materials, a shared camera, and a common metric frame.  Panel (c) reports
paired $\Delta Y$ at the highest pre-correction vertex in each occupied
1-metre X/Z cell.  Stored RGB is omitted to expose the vertical change; the
values describe implemented behaviour rather than surveyed elevation error.}
\label{fig:dsm-scene-nw-2x3}
\end{figure}

A separate representative paired building remains a single connected
component with finite vertices, valid indices, and no detected zero-area,
duplicate, or non-manifold faces. Its mean paired vertical displacement is approximately
\SI{58.0}{\metre}.  This building-level case establishes the presence of open
boundaries; the population analysis in Section~\ref{sec:ondemand-results}
evaluates their prevalence.

\section{Static ChordAtlas integration}\label{sec:static-results}

The static path converted a high-resolution source mesh into a structurally
valid MiniMesh workspace with the required indices and generator structure.
The mesh's rectangular X/Z extent overlaps approximately
\SI{0.0096}{\kilo\metre\squared} of a \SI{1.00}{\kilo\metre\squared} target,
or 0.96 per
cent.  This is an axis-aligned extent proxy, not surface area, so it can include
empty regions.  Footprint-associated assets were exported, with overlapping
OSM filter categories reported separately.

The evaluated result demonstrates

\begin{center}
cached meshes $\rightarrow$ coordinate/footprint processing $\rightarrow$
MiniMesh conversion $\rightarrow$ structurally valid workspace,
\end{center}

using retained inputs.  It establishes coordinate, footprint, format, and
workspace integration over the available geographic extent.

\section{On-demand ChordAtlas integration}\label{sec:ondemand-results}

\subsection{Independent-tile publication population and asset quality}

The tile-based on-demand set contains ten publication records with
building-level metrics and one earlier record without them.
Table~\ref{tab:on-demand-cases} aggregates these records by corrected-scene
decision.

\begin{table}[H]
\centering
\caption{Independent-tile publication records by corrected-scene decision.}
\label{tab:on-demand-cases}
\small
\begin{tabularx}{\textwidth}{XrrrX}
\toprule
Record category & Publications & Tile references & Buildings & Publication evidence\\
\midrule
Fresh correction & 5 & 90 & 11 & Modular building assets published\\
Reused correction & 5 & 91 & 8 & Modular building assets published\\
Earlier schema & 1 & 9 & --- & Publication boundary reached; no current metrics\\
\bottomrule
\end{tabularx}
\end{table}

Across the ten metric-bearing independent-tile records, 181 tile references produced 19
building assets: eleven following fresh correction and eight following cache
reuse.  Figure~\ref{fig:on-demand-publications-update} distinguishes selection
size from the number of modular outputs.

\begin{figure}[H]
\centering
\includegraphics[width=\textwidth]{on_demand_publications_update.pdf}
\caption[On-demand publication records and corrected-scene reuse.]{On-demand publication records, tile counts, building counts, and
corrected-scene reuse decisions.  The Python publication boundary precedes
Java conversion.}
\label{fig:on-demand-publications-update}
\end{figure}

Topology is the principal quality constraint in this published population.

\begin{figure}[H]
\centering
\includegraphics[width=\textwidth]{on_demand_building_topology_update.pdf}
\caption{Boundary-edge metrics for all 19 independent-tile building publications and
non-manifold-edge metrics for the ten v2 publications that record this field.}
\label{fig:on-demand-topology-update}
\end{figure}

RGB attributes are preserved throughout the independent-tile publication set:
every vertex in all 19 cropped OBJ assets carries colour data.  The
representative output in Figure~\ref{fig:on-demand-vertex-colour} shows this
appearance together with the open-edge topology identified above.

All 19 independent-tile buildings also contain MiniMesh indices.  A real-asset Java check parsed
vertex colours and constructed the generator layer successfully, confirming
appearance transfer through OBJ parsing and generator construction.

\begin{figure}[H]
\centering
\includegraphics[width=\textwidth]{on_demand_vertex_colour_building.pdf}
\caption{Oblique and plan views of a representative building-level
vertex-colour OBJ.  Colour visualises the published appearance attributes.}
\label{fig:on-demand-vertex-colour}
\end{figure}


\subsection{Spatially matched appearance references}

The independent-tile bridge extends building publication with two forms of
appearance input.  Four building assets have ready satellite roof references generated
from the tiles bound to their selections.  Each reference covers its source
footprint completely, records a footprint-derived style mask, and leaves the
published geometry unchanged.  This creates a spatially matched roof-appearance
input for the ChordAtlas style interface.

A selected-block panorama run also completed its planning, sample approval,
batch acquisition, and promotion stages.  Eight unique Street View panoramas
were acquired and indexed in the ChordAtlas workspace using the project's
local coordinate frame.  The result connects a chosen OSM footprint to a
deduplicated, selection-specific set of street-level observations.

Together, these results extend the bridge from geometry delivery to registered
appearance inputs.  For this reference-registration experiment, the
roof-reference style contract and panorama services pass their component
checks; the evaluated outcome is reference preparation and spatial
registration rather than a texture-quality measure.

\subsection{Failure containment}

Input validation rejected a duplicate footprint identity and an incomplete
imagery set before inference.  Both cases reached explicit terminal states,
demonstrating \gls{failclosed} behaviour at the request and input-validation
boundaries.  Append-only selection logs preserve these transitions.


\section{Matched London building comparison}
\label{sec:matched-building-comparison}

\subsection{Directly matched Sat3DGen routes}

The large-image and independent-tile manifests request the same target bounds
and the same two OSM footprints.  Figure~\ref{fig:matched-sat3dgen-building}
compares the main building, OSM way~369245408, with identical cameras and
metric display scale.  Both outputs store per-vertex RGB produced within their
respective Sat3DGen routes, and both
were processed against the same one-metre DSM with complete recorded vertex
coverage.  Table~\ref{tab:matched-building-metrics} reports corresponding
geometry, topology, relief, and colour-variation descriptors.

\begin{figure}[!htbp]
\centering
\includegraphics[width=\textwidth]{matched_sat3dgen_building_comparison.pdf}
\caption[Matched Sat3DGen building comparison.]{Directly matched building
comparison for OSM way~369245408.  Panels (a,c) show the fractional-feather
large-image publication and panels (b,d) the nine-mesh independent-tile
publication.  The latter retains additional sheets, spikes, and detached
fragments in this sample.  Colours are per-vertex RGB retained from the two
Sat3DGen routes, not satellite pixels projected directly onto the mesh.}
\label{fig:matched-sat3dgen-building}
\end{figure}

\begin{table}[H]
\centering
\caption{Descriptive metrics for the directly matched main-building outputs.}
\label{tab:matched-building-metrics}
\small
\begin{tabularx}{\textwidth}{p{0.18\textwidth}p{0.18\textwidth}p{0.10\textwidth}p{0.18\textwidth}p{0.12\textwidth}X}
\toprule
Route & Vertices / faces & Relief (m) & Boundary / non-manifold & Luminance s.d. & Mean edge RGB $\Delta$\\
\midrule
Large-image & 16,644 / 32,597 & 28.12 & 681 / 0 & 0.154 & 0.053\\
Independent tiles & 21,555 / 42,524 & 38.07 & 1,074 / 384 & 0.213 & 0.082\\
\bottomrule
\end{tabularx}
\end{table}

The large-image sample has a more compact surface and fewer recorded boundary
and non-manifold edges.  The independent-tile sample contains more geometry
and stronger local RGB variation.  The latter observation is consistent with
visibly higher-frequency colour detail, but it is not evidence of more accurate
colour.  Similarly, complete DSM coverage shows that both outputs were
vertically conditioned.

\subsection{Procedural appearance representations}

Figure~\ref{fig:procedural-representation-examples} places the CityEngine export
and retained ChordAtlas artefacts under a common presentation camera and metric
scale.  These are illustrative representation examples rather than a
common reconstruction benchmark.  The route-labelled ChordAtlas envelopes are
geometrically consistent with the main footprint, but their OBJ files do not
embed its OSM identifier.  The standalone ChordAtlas model has neither a
matched location nor a retained run manifest.

\begin{figure}[!htbp]
\centering
\includegraphics[width=\textwidth]{procedural_representation_examples.pdf}
\caption[Procedural model representation examples.]{Illustrative procedural
representation examples, centred without rescaling.  The CityEngine export is
a compact UV-mapped envelope with repeated facade atlases.  The two
route-labelled ChordAtlas exports use generated UV materials rather than
vertex-RGB encoding and preserve the broad building envelope.  The standalone
ChordAtlas asset is included only as a stored
contextual reference.  Different or undocumented inputs prevent a fidelity
ranking across panels.}
\label{fig:procedural-representation-examples}
\end{figure}

The CityEngine 2026 log records one initial shape exported as 120 vertex
records, 49 polygon records, 84 texture coordinates, and three materials.  Its
facades are regular, and the visible doors and windows reside in repeated
texture atlases rather than separate door or window geometry.  This supports a
conclusion about export compactness and material organisation, not about
site-specific texture accuracy \cite{esriCityEngineIntro2026,
esriCityEngineSDK2026}.

The two ChordAtlas exports contain no vertex RGB.  Instead, the large-image
route stores 27 and the independent-tile route 47 complete diffuse, normal, and
specular map sets.  Window-labelled materials and the aligned map triplets
provide a multi-channel material representation rather than a single RGB value
at each vertex.  Figure~\ref{fig:appearance-map-encoding} illustrates that
distinction.  The tile-derived export also contains one material named
\texttt{error} on ten face records.

\begin{figure}[!htbp]
\centering
\includegraphics[width=\textwidth]{appearance_map_encoding_examples.pdf}
\caption[Facade appearance-map encodings.]{Selected facade material encodings.
The CityEngine facade material uses a repeated diffuse atlas; no normal or
specular map is stored for that facade material.  The retained ChordAtlas
procedural appearance exports store aligned diffuse, normal, and specular maps
for window-labelled regions \cite{kelly2018frankengan,kellyChordAtlas2024}.  These
maps can depict frames, panes, and material response.}
\label{fig:appearance-map-encoding}
\end{figure}

The retained route-labelled artefacts place satellite-route vertex colour and
structured procedural appearance maps side by side, although the OBJ files do
not provenance-lock a direct conversion between them.
They do not show that ChordAtlas preserves the Sat3DGen triangulation,
reconstructs the true facade texture, or produces a more accurate building.

\FloatBarrier

\section{Results synthesis}

All component tests passed within their specified scope, supporting the
coordinate, cache, geometry-processing, and publication contracts across
Python and Java.  Appendix~\ref{app:evidence-record} provides the detailed test
categories and results.  Tables~\ref{tab:results-synthesis} and
\ref{tab:results-synthesis-reliability} consolidate the principal findings by
evaluation dimension.

\begin{table}[H]
\centering
\caption{Synthesis of principal findings and their evidential scope.}
\label{tab:results-synthesis}
\small
\begin{tabularx}{\textwidth}{X p{0.20\textwidth}X}
\toprule
Evaluation dimension & Evidence basis & Finding and interpretation\\
\midrule
Static ChordAtlas conversion & Structural inspection & A cached high-resolution mesh produced a valid MiniMesh workspace.\\
London data integrity & Hash and spatial analysis & Exact panorama duplication, split leakage, candidate-offset errors, and OSM support-node contamination require correction before learning use.\\
Large-image reconstruction & Matched run and round-trip inspection & Feathering reduced window-boundary discontinuities; every global-mesh vertex received colour without changing geometry, and the current ChordAtlas path retained RGB through publication.\\
Matched building representation & Shared OSM request, OBJ/MTL audit, and common-camera rendering & The large-image sample is more compact, while the tile sample has stronger colour variation and more fragments; CityEngine and ChordAtlas differences support representation-level, not accuracy, conclusions.\\
Building and appearance publication & Execution and component tests & Nineteen independent-tile assets and two current large-image assets were published; the paths retain RGB and provide either registered appearance references or the second-pass colour renderer.\\
\bottomrule
\end{tabularx}
\end{table}

\begin{table}[H]
\centering
\caption{Synthesis of coordinate and geometry-reliability findings.}
\label{tab:results-synthesis-reliability}
\small
\begin{tabularx}{\textwidth}{X p{0.20\textwidth}X}
\toprule
Evaluation dimension & Evidence basis & Finding and interpretation\\
\midrule
Coordinate control & Calculation and projection comparison & The local transform is reversible, with up to approximately \SI{19.1}{\metre} disagreement relative to EPSG:27700 displacement.\\
Stitching rule & Controlled regression & A \SI{10}{\metre}-separated fixture collapsed from six to three vertices because vertical separation was omitted.\\
DSM correction rule & Artefact comparison & Every scene vertex moved positively.\\
Building topology & Population inspection & All 19 independent-tile publications and both current large-image publications retain boundary edges and none is watertight; nine of the ten v2 independent-tile assets reporting non-manifold incidence contain such edges.\\
\bottomrule
\end{tabularx}
\end{table}
